%  SSC0903 - Computacao de Alto Desempenho
%  TB2 - Simulacao paralela em C/MPI/OpenMP da propagacao direcional de incendio
%
%  Compilar com:  pdflatex main.tex  (duas vezes, por causa das refs)
%
%  As tabelas, as figuras e os numeros citados no texto vem da pasta plot/,
%  gerada por testes/scripts/process_results.py. No Overleaf, envie este
%  arquivo junto com a pasta plot/ inteira.
\documentclass[11pt,a4paper]{article}

\usepackage[utf8]{inputenc}
\usepackage[T1]{fontenc}
\usepackage[brazilian]{babel}
\usepackage[margin=2.5cm]{geometry}
\usepackage{amsmath,amssymb}
\usepackage{graphicx}
\usepackage{booktabs}
\usepackage{multirow}
\usepackage{listings}
\usepackage{xcolor}
\usepackage{pgfplots}
\usepackage{pgfplotstable}
\usepackage{caption}
\usepackage{enumitem}
\usepackage[hidelinks]{hyperref}
\usepackage{indentfirst}
\usepackage{seqsplit}
\usepackage{float}

\pgfplotsset{compat=1.17}

% realce de codigo C/MPI/OpenMP
\definecolor{cmt}{RGB}{106,115,125}
\definecolor{kw}{RGB}{0,92,197}
\definecolor{str}{RGB}{3,106,87}
\definecolor{bg}{RGB}{248,248,248}
\lstdefinestyle{c}{
  language=C,
  basicstyle=\ttfamily\footnotesize,
  keywordstyle=\color{kw}\bfseries,
  commentstyle=\color{cmt}\itshape,
  stringstyle=\color{str},
  backgroundcolor=\color{bg},
  numbers=left, numberstyle=\tiny\color{cmt},
  frame=single, framesep=4pt, rulecolor=\color{cmt!40},
  breaklines=true, showstringspaces=false,
  tabsize=2, captionpos=b,
  morekeywords={pragma,omp,parallel,for,reduction,schedule,static,dynamic,guided,
                single,master,task,barrier,simd,collapse,shared,private,firstprivate,default,
                num_threads,nowait,long}
}

\lstset{
    style=c,
    inputencoding=utf8,
    extendedchars=true,
  literate=
    {á}{{\'a}}1 {à}{{\`a}}1 {ã}{{\~a}}1 {â}{{\^a}}1
    {é}{{\'e}}1 {ê}{{\^e}}1 {í}{{\'i}}1
    {ó}{{\'o}}1 {õ}{{\~o}}1 {ô}{{\^o}}1
    {ú}{{\'u}}1 {ç}{{\c c}}1
    {Á}{{\'A}}1 {Â}{{\^A}}1 {Ã}{{\~A}}1 {É}{{\'E}}1 {Ê}{{\^E}}1
    {Í}{{\'I}}1 {Ó}{{\'O}}1 {Õ}{{\~O}}1 {Ô}{{\^O}}1
    {Ú}{{\'U}}1 {Ç}{{\c C}}1
}
\renewcommand{\lstlistingname}{Código}

\newcommand{\TODO}[1]{\textcolor{red}{\textbf{[TODO: #1]}}}

\newcommand{\pct}{\%}

% valores, tabelas e figuras gerados pelos scripts
% \Valor{chave} escreve um numero medido (definido em plot/valores.tex) ou um TODO em vermelho se ele ainda nao existe
% \InputPlot{arquivo} carrega uma tabela ou figura de plot/ ou um TODO em vermelho se o arquivo ainda nao existe
\newcommand{\DefValor}[2]{\expandafter\def\csname val@#1\endcsname{#2}}
\newcommand{\Valor}[1]{\ifcsname val@#1\endcsname\csname val@#1\endcsname\else\TODO{#1}\fi}
\newcommand{\InputPlot}[1]{\IfFileExists{plot/#1}{\input{plot/#1}}{\TODO{arquivo \texttt{plot/\detokenize{#1}} não encontrado, rode \texttt{make resultados}}}}
\IfFileExists{plot/valores.tex}{% Gerado por testes/scripts/process_results.py, não editar à mão
\DefValor{cargaRef}{grande}
\DefValor{dimensoesRef}{$2500 \times 2500$}
\DefValor{passosRef}{100}
\DefValor{tempoSeqRef}{5{,}84}
\DefValor{melhorNP}{4}
\DefValor{melhorT}{4}
\DefValor{melhorP}{16}
\DefValor{melhorTempo}{0{,}71}
\DefValor{melhorSpeedupAbs}{8{,}25}
\DefValor{melhorEficienciaAbs}{52}
\DefValor{tempoParUmRef}{5{,}31}
\DefValor{melhorSpeedupRel}{7{,}50}
\DefValor{melhorEficienciaRel}{47}
\DefValor{fatorBase}{1{,}100}
\DefValor{comparacaoBase}{mais rápida}
\DefValor{diferencaBasePct}{9{,}1}
\DefValor{fracaoSerialPct}{7{,}5}
\DefValor{speedupMaxAmdahl}{13{,}2}
\DefValor{speedupMin}{4{,}26}
\DefValor{speedupMax}{8{,}79}
\DefValor{hibridoP}{8}
\DefValor{hibridoMelhorNP}{4}
\DefValor{hibridoMelhorT}{2}
\DefValor{hibridoMelhorTempo}{0{,}76}
\DefValor{hibridoPiorNP}{2}
\DefValor{hibridoPiorT}{4}
\DefValor{hibridoPiorTempo}{1{,}37}
\DefValor{hibridoDiferencaPct}{80{,}0}
\DefValor{numCargas}{3}
\DefValor{numConfiguracoes}{9}
\DefValor{numExecucoesPar}{135}
\DefValor{resultadoValidacao}{idênticos}
\DefValor{listaNP}{1, 2, 4}
\DefValor{listaT}{1, 2, 4}
\DefValor{maiorP}{16}
\DefValor{reps}{5}
\DefValor{numNos}{4}
\DefValor{nucleosFisicosPorNo}{4}
\DefValor{nucleosFisicosTotais}{16}
\DefValor{cpusLogicasTotais}{32}
}{}

% permite quebra de linha dentro de identificadores longos em \texttt
\emergencystretch=3em
\hbadness=10000

\title{\vspace{-1.5cm}
  \large Universidade de São Paulo\\
  \large Instituto De Ciências Matemáticas e de Computação\\
  \large SSC0903, Computação de Alto Desempenho\\[0.6cm]
  \textbf{\Large Simulação paralela em C/MPI/OpenMP da propagação\\direcional de incêndio com zonas de contenção}\\[0.2cm]
  \normalsize 2\textordmasculine \, Trabalho Prático}

\author{
  Alec Campos Aoki, 15436800 \\
  Henrique Vieira Lima, 15459372 \\
  João Gabriel Pieroli da Silva, 15678578 \\
  Pedro Augusto Ferraro Paffaro, 15483380 \\
  Pedro Henrique Perez Dias, 15484075 \\
}
\date{São Carlos, \today}

\begin{document}
\maketitle

\begin{abstract}
\noindent
Este trabalho apresenta a versão paralela híbrida, em C com MPI e OpenMP (\texttt{fire\_mpi\_omp.c}), do autômato celular que simula a propagação de incêndio em uma área florestal sujeita a vento direcional constante e a zonas de contenção ativadas em passos pré-determinados, cuja versão sequencial foi desenvolvida no primeiro trabalho. O projeto do algoritmo paralelo segue a metodologia PCAM. A matriz é dividida em blocos de linhas consecutivas, um para cada processo MPI, e dentro de cada processo as linhas do bloco são divididas entre as $T$ \emph{threads}. Como o próximo estado de uma célula depende das células ao redor, cada processo guarda também uma cópia da linha mais próxima de cada bloco vizinho, e essas cópias são atualizadas a cada passo. Toda a comunicação entre processos usa primitivas coletivas do MPI: os parâmetros são enviados a todos os processos, a matriz é distribuída no início e reunida no final, os processos vizinhos trocam suas linhas de borda e os contadores de todos os blocos são somados em totais globais. Como todos os processos recebem esses totais, todos decidem parar no mesmo passo. Dentro de cada processo, uma única região paralela OpenMP é mantida durante toda a simulação e combina \texttt{omp for}, \texttt{simd}, reduções e tarefas. As tarefas são usadas na primeira e na última linha de cada bloco, que só podem ser calculadas depois que os dados dos processos vizinhos chegam. Enquanto esses dados não chegam, as \emph{threads} calculam as outras linhas do bloco. Nas \Valor{numExecucoesPar} execuções paralelas realizadas (\Valor{numCargas} cargas, NPrc em \{\Valor{listaNP}\} e $T$ em \{\Valor{listaT}\}), os resultados foram \Valor{resultadoValidacao} aos da versão sequencial em todos os campos exceto o tempo. Na carga \Valor{cargaRef} (\Valor{dimensoesRef} células, \Valor{passosRef} passos), o tempo do núcleo da simulação caiu de $\Valor{tempoSeqRef}$ segundos na versão sequencial para $\Valor{melhorTempo}$ segundos com \Valor{melhorNP} processos de \Valor{melhorT} \emph{thread}(s), um \emph{speedup} absoluto de $\Valor{melhorSpeedupAbs}\times$ e eficiência de $\Valor{melhorEficienciaAbs}\pct$ sobre os \Valor{melhorP} núcleos utilizados.
\end{abstract}

\tableofcontents
\vspace{0.5cm}

\section{Introdução}
\label{sec:intro}

Este trabalho dá continuidade ao primeiro trabalho prático da disciplina, no qual foram desenvolvidas uma versão sequencial e uma versão paralela em memória compartilhada (C/OpenMP) de uma simulação baseada em autômatos celulares para a propagação direcional de incêndios florestais. O modelo representa a área de estudo como uma matriz bidimensional discreta, na qual o estado de cada célula no instante seguinte é determinado a partir de seu estado atual, do estado de sua vizinhança de Moore (oito vizinhos), da cobertura do solo, da umidade da vegetação, da direção e intensidade de um vento constante, e da ativação programada de zonas de contenção.

O objetivo agora é uma única versão paralela híbrida, em C com MPI e OpenMP, capaz de utilizar vários nós de um \emph{cluster} (memória distribuída) e os vários núcleos de cada nó (memória compartilhada). A simulação é determinística: para uma mesma entrada, a versão híbrida deve produzir exatamente o mesmo resultado da versão sequencial do primeiro trabalho, qualquer que seja o número de processos e de \emph{threads}.

Em relação à versão OpenMP, a principal dificuldade nova é que a matriz deixa de estar em um único espaço de endereçamento. O próximo estado de uma célula depende dos seus oito vizinhos, e as células da fronteira de um processo têm vizinhos que pertencem a outro processo. Essa dependência precisa ser resolvida por troca de mensagens a cada passo, assim como o cálculo das estatísticas globais e a decisão de parada, que precisa ser a mesma em todos os processos.

O relatório está organizado da seguinte forma. A Seção~\ref{sec:base} resume o que foi reaproveitado do primeiro trabalho. A Seção~\ref{sec:pcam} apresenta o projeto do algoritmo paralelo seguindo a metodologia PCAM. A Seção~\ref{sec:impl} descreve a implementação em C/MPI/OpenMP. A Seção~\ref{sec:validacao} valida a versão paralela contra a sequencial. A Seção~\ref{sec:ambiente} caracteriza o ambiente experimental, a Seção~\ref{sec:resultados} apresenta tempos, \emph{speedup} e eficiência, a Seção~\ref{sec:analise} discute os resultados e a Seção~\ref{sec:conclusao} traz as considerações finais. O Apêndice~\ref{sec:apendice-execucao} apresenta um guia de compilação e execução e o Apêndice~\ref{sec:apendice-tabelas} reúne as tabelas completas das demais cargas.

\section{Base reaproveitada do primeiro trabalho}
\label{sec:base}

A versão híbrida mantém as regras do modelo, as estruturas de dados e as funções de cálculo do primeiro trabalho, de modo que as diferenças de resultado e de desempenho entre as versões fiquem associadas apenas à paralelização. Os pontos a seguir foram mantidos sem alteração de lógica.

\begin{itemize}
    \item \textbf{Disposição em memória.} Cada atributo das células fica em um vetor linear separado (\emph{Structure of Arrays}), indexado por $i = l \cdot C + c$: \texttt{cobertura}, \texttt{umidade}, \texttt{estado\_atual}, \texttt{tempo\_atual}, \texttt{proximo\_estado}, \texttt{proximo\_tempo} e \texttt{ativacao}. A troca entre passo atual e próximo passo é uma permutação de ponteiros, com custo $O(1)$.
    \item \textbf{Funções de cálculo.} \texttt{fator\_cobertura}, \texttt{tempo\_queima\_inicial}, \texttt{peso\_vizinho}, \texttt{potencial\_ignicao}, \texttt{percentual\_queimado} e \texttt{percentual\_protegido} são as mesmas, assim como as funções de leitura e validação da entrada.
    \item \textbf{Otimizações de simulação.} Os deslocamentos lineares até os oito vizinhos e os pesos direcionais (que dependem só do vento, constante) são pré-calculados uma vez. A ativação das zonas e a soma das contribuições dos vizinhos são escritas sem desvios condicionais, o que permite a vetorização com \texttt{simd}. As células internas, cujos oito vizinhos sempre existem, são tratadas separadamente das células de borda.
    \item \textbf{Geração da matriz.} A cobertura e a umidade são sorteadas com \texttt{rand\_r}, célula a célula, em ordem crescente de índice, exatamente como no sequencial.
\end{itemize}

A versão sequencial entregue no primeiro trabalho (\texttt{fire\_seq.c}) é usada neste relatório como referência de corretude (Seção~\ref{sec:validacao}) e como base do \emph{speedup} absoluto (Seção~\ref{sec:resultados}).

\section{Projeto do algoritmo paralelo (PCAM)}
\label{sec:pcam}

O projeto seguiu as quatro etapas da metodologia PCAM. Nas duas primeiras (particionamento e comunicação) o problema é descrito com a granularidade mais fina possível, sem considerar a máquina nem o custo das tarefas. A máquina só é considerada a partir da aglomeração. Ela é um \emph{cluster} em que cada nó tem a sua própria memória e vários núcleos que compartilham essa memória, utilizado por NPrc processos MPI com $T$ \emph{threads} cada.

\subsection{Particionamento}
\label{sec:particionamento}

A simulação tem $P$ passos sobre $L \cdot C$ células. Os passos são inerentemente sequenciais entre si, porque o passo $p+1$ só pode ser calculado a partir do resultado completo do passo $p$. Dentro de um passo, porém, o problema pode ser dividido pelos dados, com cada célula tratada de forma independente. Buscando a granularidade mais fina possível, o trabalho de um passo $p$ é dividido nas seguintes tarefas.

\begin{enumerate}[label=(\roman*)]
    \item \textbf{Ativação.} Uma tarefa por célula verifica se \texttt{ativacao[i]} é igual a $p$ e se a célula está intacta e, nesse caso, muda seu estado para contenção. São $L \cdot C$ tarefas independentes entre si, pois cada uma lê e escreve apenas a própria célula.
    \item \textbf{Contribuição de um vizinho.} Para cada célula intacta, uma tarefa por vizinho de Moore (até oito) verifica se aquele vizinho está em chamas e, se estiver, produz o seu peso direcional $P_v$; caso contrário, produz zero. As oito tarefas de uma célula são independentes entre si.
    \item \textbf{Soma das contribuições.} Os pesos de uma célula são somados par a par por tarefas de redução (sete somas organizadas em três níveis), produzindo $S$.
    \item \textbf{Próximo estado da célula.} Uma tarefa por célula calcula o próximo estado e o próximo tempo: para uma célula intacta, calcula o potencial de ignição $I$ a partir de $S$, do fator de combustível e da umidade, e o compara com \texttt{LIMIAR}; para uma célula em chamas, decrementa o tempo de queima e verifica se ele chegou a zero; para os demais estados, apenas repete o estado. Cada tarefa escreve somente na posição da própria célula nos vetores de próximo passo.
    \item \textbf{Estatísticas.} Cada tarefa de próximo estado produz também os indicadores do próximo estado da sua célula (em chamas, nova ignição, não combustível, intacta, queimada, contenção). Para cada um desses seis contadores, $L \cdot C - 1$ tarefas de redução somam os indicadores par a par, em tempo proporcional a $\log_2 (L \cdot C)$.
    \item \textbf{Fechamento do passo.} Uma tarefa atualiza o total de ignições e o pico, troca os papéis dos vetores atual e próximo e avalia a condição de parada (nenhuma célula em chamas ou $P$ passos concluídos).
\end{enumerate}

Fora do laço de passos existem ainda as tarefas de preparação e de finalização. A geração da cobertura e da umidade não pode ser particionada: cada chamada a \texttt{rand\_r} depende do estado deixado pela chamada anterior, e a ordem das chamadas é fixa, então ela é uma única tarefa sequencial. A aplicação dos focos pode ser dividida em uma tarefa por foco, e a marcação das zonas de contenção em uma tarefa por célula de cada zona. Quando duas zonas cobrem a mesma célula, vale o menor passo de ativação entre elas. A contagem inicial dos estados é uma redução como a do item~(v). O \emph{checksum} é, por definição, uma recorrência sequencial sobre a ordem linear da matriz e forma uma única tarefa, assim como a impressão do resultado.

\subsection{Comunicação}
\label{sec:comunicacao}

Com as tarefas definidas na granularidade mais fina, as necessidades de comunicação entre elas são as seguintes.

\begin{itemize}
    \item \textbf{Entre células vizinhas.} Cada tarefa de contribuição de um vizinho (item~ii) precisa do estado atual da célula vizinha. Esse estado foi calculado no passo anterior pela tarefa de próximo estado daquela célula (item~iv) e pode ter sido alterado no passo atual pela tarefa de ativação dela (item~i). O padrão é o mesmo para todas as células e não muda ao longo da simulação: cada uma envia o seu estado para as oito células ao redor. A ordem precisa ser respeitada: todas as ativações de um passo terminam antes das contribuições, e todos os próximos estados de um passo são calculados antes de o passo seguinte começar.
    \item \textbf{Dentro de uma soma.} Cada tarefa de soma recebe os resultados de duas tarefas anteriores, tanto na soma dos pesos de uma célula (item~iii) quanto na soma das estatísticas (item~v). Uma soma só pode ser feita depois que as duas parcelas estão prontas.
    \item \textbf{Para todas as tarefas.} Algumas informações precisam chegar a todas as tarefas de todas as células. A primeira é o resultado da condição de parada, que decide se o passo seguinte existe. A segunda é o número do passo atual, usado pelas tarefas de ativação. A terceira são os parâmetros que não mudam durante a simulação (\texttt{LIMIAR}, as dimensões da matriz e os pesos direcionais, que dependem só do vento), usados pelas tarefas de contribuição e de próximo estado.
    \item \textbf{Distribuição e coleta.} A tarefa que gera a matriz precisa entregar a cada célula a sua cobertura, a sua umidade, o seu estado inicial e o seu passo de ativação. No final, o estado e o tempo de todas as células precisam chegar, na ordem linear da matriz, à tarefa que calcula o \emph{checksum}.
\end{itemize}

Não existe comunicação entre tarefas de próximo estado de células diferentes dentro do mesmo passo: todas leem apenas valores do passo atual e cada uma escreve apenas na sua posição do próximo passo.

\subsection{Aglomeração}
\label{sec:aglomeracao}

Na granularidade do particionamento, a simulação teria da ordem de dezenas de tarefas por célula por passo, cada uma com uma única comparação ou soma, e o custo de criar, sincronizar e comunicar essas tarefas seria muito maior que o trabalho útil. A aglomeração aumenta a granularidade para reduzir esse custo, considerando a máquina: entre nós diferentes a comunicação é feita por mensagens e é cara, e dentro de um mesmo nó ela é apenas um acesso à memória.

\paragraph{Tarefas de uma mesma célula.} As oito tarefas de contribuição (item~ii), as sete somas dos pesos (item~iii) e o cálculo do próximo estado (item~iv) de uma célula são aglomerados em uma única tarefa executada sequencialmente. A soma dos oito pesos vira um laço de tamanho fixo, sem desvios, que é vetorizado (\texttt{simd}) em vez de paralelizado. Com isso, toda a comunicação entre níveis da soma dos pesos deixa de existir.

\paragraph{Células em blocos de linhas.} As tarefas de células vizinhas são aglomeradas em blocos contíguos de linhas inteiras da matriz, um bloco por processo. A comunicação de vizinhança entre células do mesmo bloco se torna acesso local, e apenas as células da primeira e da última linha de cada bloco continuam dependendo de células de outro bloco. Essa dependência é resolvida guardando, em cada bloco, uma cópia da linha mais próxima de cada bloco vizinho. Essas cópias são chamadas neste relatório de linhas fantasmas e são atualizadas uma vez por passo. A escolha de blocos de linhas inteiras, e não de retângulos menores que dividissem também as colunas, se justifica por três motivos: (a) como a matriz é armazenada linearmente por linhas, um bloco de linhas e cada linha de fronteira são regiões contíguas de memória, que podem ser distribuídas e trocadas sem empacotamento; (b) cada bloco tem no máximo dois vizinhos, enquanto com retângulos seriam até oito, porque os vizinhos diagonais também contam; (c) o volume trocado por bloco e por passo é de $2C$ células, contra $(L/\text{NPrc}) \cdot C$ células processadas, uma razão de $2\,\text{NPrc}/L$ que é pequena para as dimensões de interesse. A divisão em retângulos só reduziria o volume de dados trocado para um número de processos bem maior do que o disponível.

\paragraph{Dados trocados entre blocos vizinhos.} A contribuição de um vizinho depende apenas de ele estar ou não em chamas. Por isso, só o vetor de estados precisa de linhas fantasmas; cobertura, umidade, tempo de queima e mapa de contenção nunca são trocados durante a simulação. Além disso, a ativação das zonas só transforma células intactas em contenção e nunca altera quais células estão em chamas, de modo que uma única troca por passo é suficiente.

\paragraph{Sub-blocos dentro do bloco.} Dentro de um processo, as linhas do bloco são subdivididas em $T$ sub-blocos contíguos de linhas, um por \emph{thread}. Como as \emph{threads} compartilham a memória, não há mensagens entre sub-blocos nem linhas fantasmas internas. Com sub-blocos contíguos, cada \emph{thread} escreve em uma região própria dos vetores de próximo passo. Duas \emph{threads} só escrevem em posições de memória próximas uma da outra no limite entre dois sub-blocos, o que evita que uma atrapalhe a \emph{cache} da outra.

\paragraph{Linhas de fronteira separadas das linhas internas.} A primeira e a última linha de um bloco, chamadas aqui de linhas de fronteira, são as únicas que dependem das linhas fantasmas e, portanto, da comunicação do passo. As demais são chamadas aqui de linhas internas. As duas linhas de fronteira são mantidas como duas tarefas separadas das linhas internas. Assim, as linhas internas podem ser calculadas enquanto as mensagens ainda estão em trânsito, e o processo não fica parado esperando a comunicação.

\paragraph{Estatísticas e condição de parada.} A árvore de redução sobre $L \cdot C$ células é substituída por três níveis: cada \emph{thread} conta, em variáveis privadas, os seis totais do passo no seu sub-bloco (células em chamas, novas ignições, não combustíveis, intactas, queimadas e em contenção); os totais das \emph{threads} são somados dentro do processo; e uma única redução entre os NPrc processos soma os seis totais de todos os processos de uma vez. Como o resultado dessa redução é entregue a todos os processos, cada um avalia localmente a condição de parada com os mesmos valores globais. Assim, a condição de parada não precisa de uma mensagem própria.

\paragraph{Preparação e finalização.} A leitura, a validação, a geração da matriz, a aplicação dos focos, a construção do mapa de contenção e a contagem inicial são aglomeradas em uma única tarefa de preparação, porque a geração é sequencial por natureza e as demais dependem da matriz completa (por exemplo, um foco só é válido se cair em célula combustível). O \emph{checksum} e a impressão formam a tarefa de finalização, que recebe os blocos de volta.

\paragraph{Informações calculadas em cada processo em vez de enviadas.} Duas informações de que todos os processos precisam não são transmitidas: a divisão das linhas da matriz entre os processos e os pesos direcionais dos oito vizinhos. As duas dependem apenas de valores que todos os processos já conhecem ($L$, NPrc e a configuração do vento). Por isso, cada processo faz essas contas por conta própria e chega ao mesmo resultado dos demais. Repetir uma conta tão barata em todos os processos custa menos do que enviar o resultado dela por mensagem.

O número de blocos e o de sub-blocos são parâmetros (NPrc e $T$), o que mantém a flexibilidade para diferentes quantidades de nós e de núcleos.

\subsection{Mapeamento}
\label{sec:mapeamento}

Cada bloco de linhas é mapeado em um processo MPI, e os processos são alocados aos nós do \emph{cluster} pelo \texttt{mpirun}, de acordo com o arquivo de \emph{hosts}. A quantidade de processos é escolhida em função da quantidade de nós, e a quantidade de \emph{threads} de cada processo em função da quantidade de núcleos do nó. Nenhuma das duas depende da carga de trabalho. O mapeamento é estático: a divisão é feita uma vez, antes da simulação, e não muda.

\begin{itemize}
    \item \textbf{Qual processo fica com qual bloco.} Os blocos são atribuídos em ordem de \emph{rank}: o processo $r$ recebe o $r$-ésimo bloco de linhas. Cada bloco tem $\lfloor L/\text{NPrc} \rfloor$ linhas, e os $L \bmod \text{NPrc}$ primeiros processos recebem uma linha a mais, de modo que os tamanhos diferem em no máximo uma linha. Os vizinhos de um bloco são os processos de \emph{rank} $r-1$ e $r+1$. O primeiro e o último processo têm vizinho de um lado só. No caso normal, em que há mais linhas do que processos, todos os processos recebem pelo menos $\lfloor L/\text{NPrc} \rfloor$ linhas. Se houver mais processos do que linhas, os processos excedentes ficam sem bloco e não participam da simulação.
    \item \textbf{Processo raiz.} O processo de \emph{rank} 0 executa a tarefa de preparação, distribui os blocos, recebe-os de volta ao final, calcula o \emph{checksum} e imprime o resultado. Ele também é responsável por um bloco, como os demais.
    \item \textbf{Comunicação entre processos.} As necessidades de comunicação levantadas na Seção~\ref{sec:comunicacao} são atendidas com primitivas coletivas do MPI. Os parâmetros chegam a todos os processos por \emph{broadcast}. Os blocos são distribuídos pela raiz com \emph{scatter} e reunidos ao final com \emph{gather}. A cada passo, cada processo troca as linhas de fronteira com os dois processos vizinhos, e uma redução entrega a todos os processos a soma dos seis totais do passo.
    \item \textbf{\emph{Threads}.} Os sub-blocos de um processo são mapeados nas suas $T$ \emph{threads}, que o sistema operacional distribui pelos núcleos do nó. A divisão das linhas entre as \emph{threads} é estática por padrão (\texttt{schedule(static)}), coerente com os sub-blocos contíguos definidos na aglomeração. As duas tarefas de fronteira são executadas por qualquer \emph{thread} que esteja livre. Todas as chamadas MPI de um processo são feitas pela sua \emph{thread} principal, o que corresponde ao nível de suporte a \emph{threads} \texttt{MPI\_THREAD\_FUNNELED} do MPI.
    \item \textbf{Balanceamento.} Os blocos têm praticamente o mesmo número de células, mas o custo de uma célula depende do seu estado: uma célula intacta examina os oito vizinhos, enquanto as demais apenas repetem ou decrementam um valor. Como a maior parte das células de qualquer bloco é intacta ou não combustível durante toda a simulação, e ambas são percorridas em todos os passos, a diferença de carga entre blocos é limitada. Um balanceamento dinâmico exigiria migrar linhas entre processos a cada mudança da frente de fogo, com custo de comunicação maior que o desbalanceamento que ele corrigiria, e por isso não foi adotado.
    \item \textbf{Processos e \emph{threads} por nó.} A alocação é feita em função da máquina, e não da carga de trabalho: os processos são distribuídos entre os nós do \emph{cluster} e as \emph{threads} de cada processo entre os núcleos do seu nó. O mesmo total de núcleos pode ser ocupado por poucos processos com muitas \emph{threads} ou por muitos processos com poucas \emph{threads}. A primeira opção reduz o número de mensagens e de linhas fantasmas; a segunda reduz o custo de sincronização entre \emph{threads} e mantém cada bloco menor na memória. A distribuição pretendida é a de um processo por nó, com uma \emph{thread} para cada núcleo físico do nó: ela usa todos os núcleos disponíveis sem que duas \emph{threads} disputem o mesmo núcleo, e mantém entre processos apenas a comunicação que é inevitável entre nós diferentes. Essa escolha é avaliada experimentalmente na Seção~\ref{sec:hibrido}.
\end{itemize}

A Figura~\ref{fig:decomposicao} ilustra o resultado do projeto para três processos.

\begin{figure}[H]
\centering
\begin{tikzpicture}[x=1cm, y=1cm, font=\footnotesize]
  % matriz global
  \draw[thick] (0,0) rectangle (3.2,4.5);
  \fill[blue!12]   (0,3.0) rectangle (3.2,4.5);
  \fill[orange!15] (0,1.5) rectangle (3.2,3.0);
  \fill[green!12]  (0,0)   rectangle (3.2,1.5);
  \draw[thick] (0,0) rectangle (3.2,4.5);
  \draw (0,3.0) -- (3.2,3.0);
  \draw (0,1.5) -- (3.2,1.5);
  \node at (1.6,3.75) {bloco do processo 0};
  \node at (1.6,2.25) {bloco do processo 1};
  \node at (1.6,0.75) {bloco do processo 2};
  \node[above] at (1.6,4.5) {Matriz global ($L \times C$)};

  % grade local do processo 1
  \begin{scope}[shift={(6.2,0.6)}]
    \fill[gray!25]   (0,2.9) rectangle (4.6,3.3);
    \fill[orange!35] (0,2.5) rectangle (4.6,2.9);
    \fill[orange!15] (0,0.8) rectangle (4.6,2.5);
    \fill[orange!35] (0,0.4) rectangle (4.6,0.8);
    \fill[gray!25]   (0,0.0) rectangle (4.6,0.4);
    \draw[thick] (0,0) rectangle (4.6,3.3);
    \draw (0,2.9) -- (4.6,2.9);
    \draw (0,2.5) -- (4.6,2.5);
    \draw (0,0.8) -- (4.6,0.8);
    \draw (0,0.4) -- (4.6,0.4);
    \draw[dashed] (0,1.65) -- (4.6,1.65);
    \node at (2.3,3.1)  {linha fantasma superior};
    \node at (2.3,2.7)  {linha de fronteira (tarefa)};
    \node at (2.3,2.08) {linhas internas, \emph{thread} 0};
    \node at (2.3,1.22) {linhas internas, \emph{thread} 1};
    \node at (2.3,0.6)  {linha de fronteira (tarefa)};
    \node at (2.3,0.2)  {linha fantasma inferior};
    \node[above] at (2.3,3.3) {Grade local do processo 1 ($T=2$)};
    \draw[<->, thick] (4.8,3.1) -- (5.6,3.1) node[right] {processo 0};
    \draw[<->, thick] (4.8,0.2) -- (5.6,0.2) node[right] {processo 2};
  \end{scope}
  \draw[->, thick, gray] (3.3,2.25) -- (6.1,2.25);
\end{tikzpicture}
\caption{Divisão da matriz resultante do projeto PCAM. À esquerda, os blocos de linhas de cada processo. À direita, o que o processo 1 guarda: as linhas fantasmas (cópias das linhas dos vizinhos), as duas linhas de fronteira, calculadas por tarefas, e as linhas internas, divididas entre as \emph{threads}.}
\label{fig:decomposicao}
\end{figure}

\section{Implementação em C/MPI/OpenMP}
\label{sec:impl}

O programa \texttt{fire\_mpi\_omp.c} é executado como
\begin{lstlisting}[language=bash,numbers=none]
mpirun --hostfile <arq_hosts> -np <NPrc> fire_mpi_omp <arq_entrada>
\end{lstlisting}
e segue as fases da Tabela~\ref{tab:fases}, que também indica as primitivas MPI usadas e o que entra no trecho cronometrado.

\begin{table}[H]
\centering
\small
\begin{tabular}{@{}lp{0.36\textwidth}lc@{}}
\toprule
Fase & O que acontece & Primitiva MPI & Cronometrada \\
\midrule
Preparação      & A raiz lê, valida, gera a matriz, aplica focos, monta o mapa de contenção e conta os estados & -- & Não \\
Configuração    & Status da leitura, parâmetros e contagens iniciais são transmitidos para todos & \texttt{MPI\_Bcast} & Não \\
Vizinhança      & Processos com linhas são organizados em fila, cada um com o vizinho de cima e o de baixo & \texttt{MPI\_Comm\_split}, \texttt{MPI\_Cart\_create} & Não \\
Distribuição    & Cada processo recebe o seu bloco de cada vetor & \texttt{MPI\_Scatterv} & Não \\
Simulação       & Ativação, troca de fronteiras, próximo estado, estatísticas globais, troca dos vetores e parada & \texttt{MPI\_Neighbor\_alltoallv}, \texttt{MPI\_Allreduce} & Sim \\
Junção          & A raiz recebe os blocos de estado e de tempo & \texttt{MPI\_Gatherv} & Não \\
Finalização     & A raiz calcula \emph{checksum} e percentuais e imprime & -- & Não \\
\bottomrule
\end{tabular}
\caption{Fases da execução, primitivas MPI utilizadas e relação com o trecho cronometrado.}
\label{tab:fases}
\end{table}

Toda a comunicação usa primitivas coletivas, isto é, chamadas em que todos os processos do grupo participam ao mesmo tempo. Não há \texttt{MPI\_Send} ou \texttt{MPI\_Recv} no programa.

As mensagens trocadas a cada passo (linhas de fronteira e redução dos totais) estão dentro do trecho cronometrado, porque fazem parte da simulação. Ficam fora dele apenas a distribuição inicial da matriz e a junção final. A primeira faz parte da preparação dos dados, como a geração da matriz, e a segunda serve ao cálculo do \emph{checksum}, e o enunciado exclui as duas coisas do tempo medido.

\subsection{Organização do código}

O fluxo principal não fica concentrado em \texttt{main}. Cada fase da execução tem uma função própria (\texttt{compartilhar\_configuracao}, \texttt{decompor\_dominio}, \texttt{distribuir\_grade}, \texttt{simular}, \texttt{juntar\_grade} e \texttt{imprimir\_resultado}), e \texttt{main} apenas encadeia essas chamadas e trata os dois encerramentos antecipados (entrada inválida e processo sem linhas). Os dados da divisão da matriz (quantas linhas cada processo tem, quem são os vizinhos e os comunicadores MPI) ficam na estrutura \texttt{DECOMPOSICAO}, e os deslocamentos e pesos dos oito vizinhos de uma célula na estrutura \texttt{VIZINHANCA}. A região paralela persistente e todas as diretivas OpenMP do laço ficam juntas em \texttt{simular}.

\subsection{Preparação e envio da configuração}

Apenas o processo raiz abre o arquivo de entrada. Ele executa as mesmas funções de leitura e validação do primeiro trabalho e gera a matriz global. Em seguida, a raiz envia a todos os processos, com um único \emph{broadcast}, um vetor de inteiros com os parâmetros da simulação e as contagens iniciais dos estados (Código~\ref{lst:bcast}). A primeira posição desse vetor indica se a leitura deu certo. Se a entrada for inválida, a raiz imprime a mensagem de erro e todos os processos, avisados por essa posição, encerram sem imprimir resultado.

\begin{lstlisting}[caption={Envio do resultado da leitura, dos parâmetros e das contagens iniciais a todos os processos.},label={lst:bcast}]
/* em compartilhar_configuracao() */
if (rank == PROC_RAIZ)
    config[CFG_STATUS] = preparar_simulacao(argc, argv, globais, config);

MPI_Bcast(config, CFG_TAMANHO, MPI_INT, PROC_RAIZ, MPI_COMM_WORLD);
return config[CFG_STATUS] == LEITURA_OK;

/* em main() */
if (!compartilhar_configuracao(argc, argv, suporte_threads, rank, &globais, config)) {
    liberar_cels(&globais);
    MPI_Finalize();
    return EXIT_FAILURE;
}
\end{lstlisting}

A geração da cobertura e da umidade permanece sequencial na raiz. Essa é a única forma de reproduzir a sequência de \texttt{rand\_r} exigida pelo enunciado, e ela fica fora do trecho cronometrado.

\subsection{Divisão da matriz entre os processos}

Todos os processos calculam a mesma divisão das linhas (\texttt{dividir\_linhas}), de modo que ela não precisa ser transmitida. Os processos que receberam ao menos uma linha são reunidos em um comunicador próprio. Sobre ele é criada uma topologia cartesiana de uma dimensão, recurso do MPI que organiza os processos em fila, na ordem dos blocos, e informa a cada um quem é o vizinho de cima e o de baixo. Por fim, um tipo de dado MPI representa uma linha inteira da matriz, para que os envios e recebimentos sejam contados em linhas (Código~\ref{lst:topologia}).

\begin{lstlisting}[caption={Comunicador dos processos com linhas, fila de processos e tipo de linha, em \texttt{decompor\_dominio} (prefixo \texttt{dec->} dos campos omitido).},label={lst:topologia}]
MPI_Comm_split(MPI_COMM_WORLD, linhas_locais > 0 ? 1 : MPI_UNDEFINED, rank, &comm_ativos);

int num_ativos = num_processos < L ? num_processos : L;
int periodico  = 0;
MPI_Cart_create(comm_ativos, 1, &num_ativos, &periodico, 0, &comm_grade);
MPI_Cart_shift(comm_grade, 0, 1, &vizinho_cima, &vizinho_baixo);

MPI_Type_contiguous(C, MPI_INT, &tipo_linha);
MPI_Type_commit(&tipo_linha);
\end{lstlisting}

Cada processo aloca os sete vetores do seu bloco com duas linhas a mais, a fantasma superior (linha local 0) e a fantasma inferior (linha local \texttt{linhas\_locais + 1}). As linhas reais ocupam as posições 1 a \texttt{linhas\_locais}. As linhas fantasmas dos dois vetores de estado são iniciadas como não combustíveis. Na primeira e na última linha da matriz, onde não há processo vizinho, elas nunca são sobrescritas e por isso nunca contribuem para o potencial de ignição. Com isso, a regra de ignorar vizinhos fora da matriz é atendida sem nenhum teste de limite de linha. Somente os limites de coluna continuam sendo verificados, nas células da primeira e da última coluna.

A raiz distribui os blocos de \texttt{cobertura}, \texttt{umidade}, \texttt{estado\_atual}, \texttt{tempo\_atual} e \texttt{ativacao} com cinco chamadas a \texttt{MPI\_Scatterv}, com as quantidades e os deslocamentos expressos em linhas.

\subsection{Laço da simulação}
\label{sec:laco}

A simulação de cada processo ocorre dentro de uma única região paralela OpenMP, aberta antes do laço de passos e fechada depois dele (Código~\ref{lst:regiao}). Cada passo tem quatro etapas, numeradas no código.

\begin{lstlisting}[caption={Estrutura da região paralela persistente de \texttt{fire\_mpi\_omp.c} (lista de variáveis compartilhadas omitida).},label={lst:regiao}]
#pragma omp parallel num_threads(T) default(none) shared(...)
{
while (passo_atual < P && cnt->em_chamas > 0) {
    int *estado_atual   = celulas->estado_atual;  /* copias privadas dos ponteiros */
    const int *ativacao = celulas->ativacao;

    #pragma omp for simd schedule(SCHED)          /* (1) ativacao das zonas */
    for (long long i = primeira_celula_real; i < fim_celulas_reais; i++) {
        int deve_ativar = (ativacao[i] == passo_atual) & (estado_atual[i] == ESTADO_INTACTA);
        estado_atual[i] = deve_ativar ? ESTADO_CONTENCAO : estado_atual[i];
    }

    #pragma omp master                            /* (2) fronteiras */
    {
        trocar_fronteiras(celulas->estado_atual, dec, C);

        #pragma omp task default(none) shared(borda_superior, celulas, ...)
        borda_superior = atualizar_linha(celulas, 1, C, LIMIAR, ...);

        if (linhas_locais > 1) {
            #pragma omp task default(none) shared(borda_inferior, celulas, ...)
            borda_inferior = atualizar_linha(celulas, linhas_locais, C, LIMIAR, ...);
        }
    }

    #pragma omp for schedule(SCHED) reduction(+ : proximo_celulas_em_chamas, \
        ignicoes_no_passo, proximo_nao_combustiveis, proximo_intactas,       \
        proximo_queimadas, proximo_contencao)     /* (3) linhas internas */
    for (int l = 2; l < linhas_locais; l++) {
        ESTATISTICAS linha = atualizar_linha(celulas, l, C, LIMIAR, ...);
        proximo_celulas_em_chamas += linha.em_chamas;
        /* ... demais contadores ... */
    }

    #pragma omp master                            /* (4) fechamento do passo */
    {
        empacotar_estatisticas(estatisticas_locais, miolo, borda_superior, borda_inferior);
        somar_estatisticas(estatisticas_locais, estatisticas_globais, dec->comm_grade);
        trocar_buffers(celulas);
        atualizar_contadores(cnt, pico, estatisticas_globais, passo_atual);
        /* ... zera os acumuladores ... */
        passo_atual++;
    }
    #pragma omp barrier
}
}
\end{lstlisting}

\paragraph{(1) Ativação das zonas.} Um \texttt{omp for simd} percorre as células reais do bloco e aplica a ativação sem desvios condicionais, como no primeiro trabalho. A barreira implícita ao final garante que todas as ativações do passo terminaram antes de qualquer leitura do estado atual.

\paragraph{(2) Troca de fronteiras e tarefas de borda.} A \emph{thread} principal envia a primeira e a última linha real do vetor de estados e recebe as duas linhas fantasmas, com uma única chamada a \texttt{MPI\_Neighbor\_alltoallv}, que envia e recebe dos dois vizinhos da fila de uma vez (Código~\ref{lst:halo}). Como \texttt{omp master} não tem barreira implícita, as demais \emph{threads} seguem diretamente para a etapa (3) e já começam a calcular as linhas internas, que não dependem das linhas fantasmas. Quando a troca termina, a \emph{thread} principal cria uma tarefa para cada linha de fronteira e também se junta à etapa (3). As tarefas são executadas por qualquer \emph{thread} que chegue à barreira da etapa (3), e essa barreira só libera as \emph{threads} depois que o laço e as duas tarefas terminaram.

\begin{lstlisting}[caption={Troca das linhas de fronteira com os dois processos vizinhos (posições contadas em linhas).},label={lst:halo}]
const int quantidades[2]     = {1, 1};
const int desloc_envio[2]    = {0, linhas_locais - 1};
const int desloc_recebido[2] = {0, linhas_locais + 1};

MPI_Neighbor_alltoallv(estado + C, quantidades, desloc_envio, dec->tipo_linha,
                       estado, quantidades, desloc_recebido, dec->tipo_linha, dec->comm_grade);
\end{lstlisting}

Nessa chamada, a ordem dos vizinhos é sempre o de cima e depois o de baixo. Quando um deles não existe, a parte correspondente da troca é ignorada pelo MPI e a linha fantasma permanece não combustível. Essa primitiva só existe a partir da versão 3 do MPI. Para versões anteriores, o código seleciona em tempo de compilação uma alternativa com duas chamadas a \texttt{MPI\_Sendrecv}.

\paragraph{(3) Linhas internas.} Um \texttt{omp for} distribui as linhas 2 a \texttt{linhas\_locais - 1} entre as \emph{threads}. Cada iteração chama \texttt{atualizar\_linha}, a mesma função usada pelas tarefas de borda, que contém a lógica de atualização do primeiro trabalho: a escolha do próximo estado conforme o estado atual, a soma vetorizada (\texttt{omp simd reduction}) dos oito vizinhos das células internas e o tratamento separado da primeira e da última coluna. A função lê apenas vetores do passo atual e escreve apenas posições da própria linha nos vetores de próximo passo, de modo que linhas diferentes podem ser calculadas simultaneamente. As estatísticas das linhas são acumuladas por \texttt{reduction}.

\paragraph{(4) Fechamento do passo.} A \emph{thread} principal soma os totais das linhas internas com os devolvidos pelas duas tarefas de borda e reduz os seis totais entre os processos com uma única chamada a \texttt{MPI\_Allreduce}. Em seguida troca os ponteiros dos vetores atual e próximo, atualiza os contadores globais, o total e o pico de ignições e incrementa o passo. A barreira explícita depois dessa etapa impede que as demais \emph{threads} reavaliem a condição do \texttt{while} antes de esses valores estarem prontos.

\subsection{Sincronização, condições de corrida e determinismo}
\label{sec:sincronizacao}

O programa não usa \texttt{critical}, \texttt{atomic} nem \emph{locks}. A ausência de condições de corrida decorre da estrutura do passo:

\begin{itemize}
    \item o estado atual só é escrito na etapa (1), por iterações independentes, e nas linhas fantasmas na etapa (2), que não são lidas no cálculo das linhas internas; as tarefas de borda, que leem as linhas fantasmas, só são criadas depois que a troca terminou;
    \item os vetores de próximo passo são escritos por linha, e cada linha pertence a exatamente uma iteração do laço ou a exatamente uma das duas tarefas OpenMP (\texttt{omp task}) de borda;
    \item os totais das linhas internas são variáveis de redução, e cada tarefa de borda escreve em uma estrutura própria, lida somente depois da barreira;
    \item a troca dos ponteiros e a atualização dos contadores e do passo são feitas por uma única \emph{thread}, entre duas barreiras.
\end{itemize}

As chamadas MPI são feitas apenas pela \emph{thread} principal, dentro de \texttt{omp master}, e por isso o programa solicita em \texttt{MPI\_Init\_thread} o nível \texttt{MPI\_THREAD\_FUNNELED}, que é o nível de suporte a \emph{threads} em que apenas a \emph{thread} principal chama o MPI.

A condição de parada é a mesma para todos em dois níveis. Entre processos, o \texttt{MPI\_Allreduce} entrega a todos o mesmo total global de células em chamas, e \texttt{passo\_atual} evolui de forma idêntica, de modo que todos decidem parar no mesmo passo. Dentro de um processo, a barreira ao final do passo garante que todas as \emph{threads} avaliem a condição com os mesmos valores.

O resultado é independente de NPrc e de $T$ porque cada célula é calculada exatamente uma vez por passo, a partir apenas de valores do passo anterior (do próprio bloco ou copiados nas linhas fantasmas), e porque todas as estatísticas são somas de inteiros, que não dependem da ordem em que são combinadas.

\subsection{Junção do resultado e medição do tempo}

Ao final, a raiz recebe os blocos de \texttt{estado\_atual} e \texttt{tempo\_atual} com \texttt{MPI\_Gatherv}, na ordem das linhas, e calcula o \emph{checksum} sequencialmente sobre a matriz global, exatamente como no sequencial. Os percentuais e a impressão também ficam na raiz.

O tempo é medido pela raiz com \texttt{omp\_get\_wtime()}. O cronômetro é disparado depois de uma \texttt{MPI\_Barrier}, que garante que todos os processos concluíram a preparação, e é parado ao fim da região paralela. O trecho cronometrado contém tudo o que o enunciado determina (ativação das zonas, propagação, atualização de estados e de tempos, estatísticas, troca das matrizes e condição de parada) e, por consequência, também a comunicação que essas etapas exigem a cada passo: a troca de fronteiras e a redução global. A distribuição inicial e a junção final não são cronometradas, por serem, respectivamente, parte da preparação dos dados e do cálculo do \emph{checksum}, que o enunciado exclui do trecho cronometrado.

\subsection{Complexidade}

O trabalho total por passo continua sendo $O(L \cdot C)$, agora dividido entre NPrc processos com $T$ \emph{threads}. A cada passo, cada processo também troca duas linhas de $C$ células e participa de uma redução entre NPrc processos. Em uma situação ideal, o tempo de um passo é
$$
O\left(\frac{L \cdot C}{\text{NPrc} \cdot T} + C + \log \text{NPrc}\right),
$$
e o da simulação é esse valor multiplicado por $P$. O primeiro termo é o trabalho útil e diminui com o aumento do paralelismo; os outros dois são o custo da comunicação e não diminuem. Para $L$ fixo, aumentar NPrc reduz o número de linhas por bloco, e a razão entre comunicação e computação, $2\,\text{NPrc}/L$, cresce.

Quanto à memória, cada processo guarda sete vetores de $(L/\text{NPrc} + 2) \cdot C$ inteiros. A raiz guarda adicionalmente a matriz global usada na geração e no \emph{checksum} (cinco vetores de $L \cdot C$ inteiros durante a preparação, dois depois da distribuição), de modo que o consumo dela permanece $O(L \cdot C)$. Esse é um limite conhecido desta implementação: o tamanho máximo do problema é determinado pela memória do nó da raiz, consequência direta de a geração e o \emph{checksum} serem sequenciais por definição.

\section{Validação}
\label{sec:validacao}

\subsection{Cargas de teste}

Foram utilizadas as três cargas do primeiro trabalho (Tabela~\ref{tab:entradas}). O valor de $T$ do arquivo de entrada é substituído em cada execução pelo número de \emph{threads} da configuração avaliada, e os demais parâmetros são mantidos.

\begin{table}[H]
\centering
\small
\begin{tabular}{@{}llrrrlrrr@{}}
\toprule
Carga & $L \times C$ & $P$ & LIMIAR & Seed & Direção & $V$ & $F$ & $Z$ \\
\midrule
Pequena & $400 \times 500$   & 80  & 35 & 2027 & leste    & 3 & 3  & 2 \\
Média   & $1200 \times 1500$ & 100 & 35 & 3031 & sudeste  & 2 & 6  & 4 \\
Grande  & $2500 \times 2500$ & 100 & 35 & 4099 & nordeste & 4 & 10 & 6 \\
\bottomrule
\end{tabular}
\caption{Cargas utilizadas na validação e nos experimentos.}
\label{tab:entradas}
\end{table}

\subsection{Metodologia}

A validação tem duas partes. A primeira é uma bateria automática de testes de corretude, que compara as onze primeiras linhas da saída de \texttt{fire\_mpi\_omp} com um gabarito produzido por \texttt{fire\_seq}, para cada combinação de NPrc e $T$, sobre as três cargas e sobre quinze casos de borda: matrizes com uma, duas e três linhas (inclusive com mais processos do que linhas), com uma e duas colunas e com uma única célula; $P=0$; nenhum foco; nenhuma ignição; fogo que se extingue antes de $P$ passos; zonas sobrepostas e zonas sobre as fronteiras entre blocos; e intensidades de vento mínima e máxima. A mesma bateria confere que 25 entradas inválidas, cobrindo todas as verificações exigidas pelo enunciado, além de chamadas com argumentos incorretos, são rejeitadas com mensagem de erro, sem impressão de resultado e com código de saída diferente de zero.

A segunda parte usa as próprias execuções de desempenho: todos os campos de saída, exceto o tempo, de todas as execuções paralelas cronometradas são comparados com os da execução sequencial da mesma carga.

\subsection{Resultados da validação}

A Tabela~\ref{tab:validacao} apresenta a saída de cada carga e o resultado da comparação. Nas \Valor{numExecucoesPar} execuções paralelas, os campos foram \Valor{resultadoValidacao} aos do sequencial.

\begin{table}[H]
\centering
\small
\resizebox{\textwidth}{!}{%
\InputPlot{tab_validacao.tex}}
\caption{Saída da versão sequencial por carga e conferência com todas as execuções da versão paralela.}
\label{tab:validacao}
\end{table}

\subsection{Independência do número de processos e de \emph{threads}}

A Tabela~\ref{tab:independencia} verifica diretamente o requisito de independência do resultado em relação a NPrc e a $T$, para a carga intermediária: \texttt{total\_ignicoes}, \texttt{pico\_ignicoes} e \emph{checksum} em cada configuração medida.

\begin{table}[H]
\centering
\small
\InputPlot{tab_independencia.tex}
\caption{Saída da versão paralela em função de NPrc e de $T$ (carga intermediária).}
\label{tab:independencia}
\end{table}

\section{Ambiente experimental}
\label{sec:ambiente}

Os experimentos foram executados no \emph{cluster} do LASDPC/ICMC-USP, nos nós \texttt{hal01} a \texttt{hal04}. As informações da Tabela~\ref{tab:ambiente} foram coletadas automaticamente no \texttt{hal01}, de onde as execuções foram disparadas, com os comandos do Código~\ref{lst:coleta}.

\begin{table}[H]
\centering
\small
\InputPlot{tab_ambiente.tex}
\caption{Configuração do ambiente de execução.}
\label{tab:ambiente}
\end{table}

\begin{lstlisting}[language=bash,caption={Coleta de informações sobre o meio de execução.},basicstyle=\ttfamily\footnotesize,label={lst:coleta}]
lscpu                 # modelo do processador, núcleos físicos e lógicos, cache
cat /etc/os-release   # sistema operacional
free -h               # memória disponível
mpicc --version       # compilador usado pelo mpicc
mpirun --version      # implementação e versão do MPI
echo | mpicc -fopenmp -dM -E - | grep -i openmp   # especificação OpenMP suportada
\end{lstlisting}

Os quatro nós utilizados têm o mesmo processador, conferido com \texttt{lscpu} em cada um deles, e são interligados por uma rede Gigabit Ethernet.

Cada configuração foi executada \Valor{reps} vezes e o valor reportado é a mediana dos tempos. O \texttt{mpirun} foi chamado com as \emph{flags} indicadas na Tabela~\ref{tab:ambiente}. A opção \texttt{--map-by node} coloca um processo em cada nó, de modo que uma execução com NPrc processos usa NPrc nós, e \texttt{--bind-to none} impede que o Open MPI prenda cada processo a um único núcleo, o que faria as $T$ \emph{threads} de um processo disputarem o mesmo núcleo.

\section{Resultados de desempenho}
\label{sec:resultados}

Em toda esta seção, $p = \text{NPrc} \cdot T$ é o total de núcleos usados por uma configuração. Como cada processo fica em um nó e $T$ não passa de quatro, que é o número de núcleos físicos de um nó, cada \emph{thread} tem um núcleo físico próprio em todas as configurações. As métricas são definidas por
$$
S_{\mathrm{abs}}(\text{NPrc}, T) = \frac{t_{\mathrm{seq}}}{t_{\mathrm{par}}(\text{NPrc}, T)},
\qquad
S_{\mathrm{rel}}(\text{NPrc}, T) = \frac{t_{\mathrm{par}}(1, 1)}{t_{\mathrm{par}}(\text{NPrc}, T)},
\qquad
E = \frac{S}{p},
$$
em que $t_{\mathrm{seq}}$ é o tempo da versão sequencial do primeiro trabalho e $t_{\mathrm{par}}(1,1)$ é o tempo da própria versão híbrida com um processo e uma \emph{thread}. O \emph{speedup} absoluto mede o ganho sobre o melhor programa sequencial disponível; o relativo mede a escalabilidade da versão paralela em si.

\subsection{Tempos de execução}
\label{sec:tempos-execucao}

A Tabela~\ref{tab:resumo} resume, para cada carga, o tempo sequencial, o tempo da versão paralela com um processo e uma \emph{thread} e a melhor configuração medida. A Figura~\ref{fig:tempos} apresenta os mesmos tempos em escala logarítmica.

\begin{table}[H]
\centering
\small
\resizebox{\textwidth}{!}{%
\InputPlot{tab_resumo.tex}}
\caption{Tempo do trecho cronometrado (mediana, em segundos), melhor configuração e \emph{speedups} de cada carga.}
\label{tab:resumo}
\end{table}

\begin{figure}[H]
\centering
\InputPlot{fig_tempos.tex}
\caption{Tempo do núcleo da simulação na versão sequencial, na versão paralela com um processo e uma \emph{thread} e na melhor configuração paralela, em escala logarítmica.}
\label{fig:tempos}
\end{figure}

A Tabela~\ref{tab:tempos} detalha o tempo da carga \Valor{cargaRef} para todas as combinações de NPrc e $T$. O menor tempo, em negrito, é de $\Valor{melhorTempo}$ segundos, com \Valor{melhorNP} processos e \Valor{melhorT} \emph{thread}(s) por processo. As tabelas equivalentes das demais cargas estão no Apêndice~\ref{sec:apendice-tabelas}.

\begin{table}[H]
\centering
\small
\InputPlot{tab_tempos_grande.tex}
\caption{Tempo do trecho cronometrado (mediana, em segundos) em função de NPrc e de $T$, carga grande.}
\label{tab:tempos}
\end{table}

\subsection{\emph{Speedup} e eficiência}
\label{sec:speedup}

As Tabelas~\ref{tab:speedup} e~\ref{tab:eficiencia} apresentam o \emph{speedup} absoluto e a eficiência absoluta de cada configuração para a carga grande, e as Figuras~\ref{fig:speedup} e~\ref{fig:eficiencia} mostram as mesmas grandezas em função do total de núcleos, com uma curva para cada valor de $T$.

\begin{table}[H]
\centering
\small
\InputPlot{tab_speedup_grande.tex}
\caption{\emph{Speedup} absoluto $S_\text{abs}$ em função de NPrc e de $T$, carga grande.}
\label{tab:speedup}
\end{table}

\begin{table}[H]
\centering
\small
\InputPlot{tab_eficiencia_grande.tex}
\caption{Eficiência absoluta $E_\text{abs} = S_\text{abs}/(\text{NPrc} \cdot T)$ em função de NPrc e de $T$, carga grande.}
\label{tab:eficiencia}
\end{table}

\begin{figure}[H]
\centering
\InputPlot{fig_speedup.tex}
\caption{\emph{Speedup} absoluto em função do total de núcleos $p$, com uma curva por número de \emph{threads} por processo, e reta de \emph{speedup} ideal.}
\label{fig:speedup}
\end{figure}

\begin{figure}[H]
\centering
\InputPlot{fig_eficiencia.tex}
\caption{Eficiência absoluta em função do total de núcleos $p$, com uma curva por número de \emph{threads} por processo.}
\label{fig:eficiencia}
\end{figure}

\subsection{Escalabilidade por total de núcleos}
\label{sec:escalabilidade}

A Tabela~\ref{tab:escalabilidade} seleciona, para cada total de núcleos $p$, a combinação de NPrc e $T$ de menor tempo na carga \Valor{cargaRef}, e apresenta para ela os \emph{speedups} absoluto e relativo, as eficiências e a métrica de Karp--Flatt,
$$
e(p) = \frac{1/S_{\mathrm{abs}} - 1/p}{1 - 1/p},
$$
que estima experimentalmente a fração do tempo que não se beneficia do aumento de $p$ (parte serial somada ao \emph{overhead} de paralelização).

\begin{table}[H]
\centering
\small
\InputPlot{tab_escalabilidade.tex}
\caption{Melhor configuração para cada total de núcleos $p$: tempo, \emph{speedup}, eficiência e métrica de Karp--Flatt.}
\label{tab:escalabilidade}
\end{table}

\subsection{Processos ou \emph{threads}: configurações com o mesmo total de núcleos}
\label{sec:hibrido}

A Tabela~\ref{tab:hibrido} agrupa as configurações que usam o mesmo total de núcleos com combinações diferentes de processos e \emph{threads}, e a Figura~\ref{fig:hibrido} compara os tempos do maior desses grupos.

\begin{table}[H]
\centering
\small
\InputPlot{tab_hibrido.tex}
\caption{Configurações com o mesmo total de núcleos $p$ e diferentes combinações de NPrc e $T$. Em negrito, a mais rápida de cada grupo.}
\label{tab:hibrido}
\end{table}

\begin{figure}[H]
\centering
\InputPlot{fig_hibrido.tex}
\caption{Tempo do núcleo da simulação para as combinações de NPrc e $T$ com o maior total de núcleos medido em mais de uma combinação.}
\label{fig:hibrido}
\end{figure}

Com $p = \Valor{hibridoP}$, a combinação mais rápida foi a de \Valor{hibridoMelhorNP} processo(s) com \Valor{hibridoMelhorT} \emph{thread}(s) ($\Valor{hibridoMelhorTempo}$\,s) e a mais lenta a de \Valor{hibridoPiorNP} processo(s) com \Valor{hibridoPiorT} \emph{thread}(s) ($\Valor{hibridoPiorTempo}$\,s), uma diferença de $\Valor{hibridoDiferencaPct}\pct$.

\section{Análise dos resultados}
\label{sec:analise}

\subsection{Versão paralela com um único núcleo}

Com um processo e uma \emph{thread}, a versão híbrida executa a carga \Valor{cargaRef} em $\Valor{tempoParUmRef}$\,s, contra $\Valor{tempoSeqRef}$\,s da versão sequencial, ou seja, é $\Valor{diferencaBasePct}\pct$ \Valor{comparacaoBase} (razão $t_\text{seq}/t_\text{par}(1,1) = \Valor{fatorBase}$). Dois efeitos atuam em sentidos opostos nessa comparação. A favor da versão híbrida, as linhas fantasmas eliminam o teste de limites de linha, e apenas as células da primeira e da última coluna passam pelo caminho com verificações, enquanto no sequencial todo o perímetro passa por ele e toda célula intacta testa os quatro limites. Contra ela, pesam as barreiras e as reduções do OpenMP, a criação das duas tarefas e as duas chamadas MPI por passo, que continuam existindo mesmo sem nenhum paralelismo. Como essa razão é próxima de 1, os \emph{speedups} absoluto e relativo ficam próximos entre si, e o \emph{speedup} absoluto pode ser lido como ganho de paralelização.

\subsection{Ganho com o aumento do número de núcleos}

Considerando a melhor configuração de cada carga (Tabela~\ref{tab:resumo}), o \emph{speedup} absoluto variou de $\Valor{speedupMin}$ a $\Valor{speedupMax}\times$. Na carga \Valor{cargaRef}, a melhor configuração usa $p = \Valor{melhorP}$ núcleos e atinge $S_\text{abs} = \Valor{melhorSpeedupAbs}$ ($E_\text{abs} = \Valor{melhorEficienciaAbs}\pct$) e $S_\text{rel} = \Valor{melhorSpeedupRel}$ ($E_\text{rel} = \Valor{melhorEficienciaRel}\pct$).

A Tabela~\ref{tab:escalabilidade} mostra que o \emph{speedup} cresce de forma quase ideal até $p = 8$ e praticamente para de crescer depois disso. O \emph{speedup} relativo é de $1{,}96$ com dois núcleos, $3{,}84$ com quatro e $6{,}96$ com oito, mas chega a apenas $7{,}50$ com dezesseis. A Tabela~\ref{tab:eficiencia} mostra de onde vem essa queda. Aumentar o número de processos quase não reduz a eficiência: com uma \emph{thread} por processo, ela vai de $1{,}10$ com um processo para $1{,}06$ com quatro. Aumentar o número de \emph{threads} tem dois comportamentos diferentes: de uma para duas \emph{threads} a eficiência se mantém (de $1{,}10$ para $1{,}01$ com um processo), e de duas para quatro ela cai pela metade (para $0{,}54$). Em tempo, dobrar as \emph{threads} de duas para quatro reduz o tempo em apenas 6\pct{} a 7\pct{}, com um, dois ou quatro processos (Tabela~\ref{tab:tempos}).

Essa perda não vem de \emph{hyperthreading}, porque nenhuma configuração usa mais \emph{threads} do que os quatro núcleos físicos de um nó. Também não vem da comunicação entre processos, que aparece ao aumentar NPrc e custa pouco. O limite está dentro de cada nó: a terceira e a quarta \emph{thread} de um processo quase não ajudam. Uma explicação compatível com as medidas é a disputa pela memória do nó. As \emph{threads} de um processo leem e escrevem os mesmos vetores e compartilham o acesso à memória, enquanto cada processo novo fica em outro nó, com memória própria. Uma restrição de núcleos imposta pelo \texttt{mpirun} foi descartada: o processo enxerga as oito CPUs lógicas do nó, e a execução com quatro \emph{threads} leva o mesmo tempo com e sem o \texttt{mpirun} ($2{,}70$\,s na carga grande). A disputa pela memória, porém, não foi isolada por um experimento específico e deve ser lida como hipótese.

\subsection{Efeito do tamanho do problema}

A razão entre comunicação e computação de um bloco é $2\,\text{NPrc}/L$ (Seção~\ref{sec:aglomeracao}): quanto menor a matriz e maior o número de processos, maior o peso relativo da troca de fronteiras e da redução global, cujo custo por passo não diminui com o tamanho do bloco. Esperava-se, portanto, que a carga pequena ($L=400$) apresentasse os menores \emph{speedups} e que as cargas maiores aproveitassem melhor o aumento de NPrc.

A Tabela~\ref{tab:resumo} confirma a previsão para a carga pequena. O melhor \emph{speedup} absoluto dela é de $4{,}26$, cerca de metade do obtido nas outras duas cargas, e com quatro processos de uma \emph{thread} ela chega a apenas $2{,}21$, contra $4{,}14$ da carga média e $4{,}23$ da grande. Os tempos da carga pequena são de dezenas de milissegundos e variaram bastante entre as repetições, de modo que a ordem entre as configurações mais rápidas dela não é confiável. Entre as cargas média e grande a previsão não se confirma: os \emph{speedups} são praticamente iguais em todas as configurações com até duas \emph{threads} por processo (por exemplo, $7{,}64$ e $7{,}65$ com quatro processos de duas \emph{threads}), e com quatro \emph{threads} a carga média fica um pouco à frente ($8{,}79$ contra $8{,}25$). A partir da carga média, portanto, o custo da comunicação entre processos já é pequeno o bastante para não diferenciar as cargas.

\subsection{Processos ou \emph{threads}}

A comparação da Seção~\ref{sec:hibrido} isola o efeito de como um mesmo total de núcleos é dividido entre processos e \emph{threads}. Mais processos significam blocos menores, mais linhas fantasmas e mais participantes na troca e na redução de cada passo, mas também menos \emph{threads} disputando as barreiras de uma mesma região paralela e, quando os processos estão em nós diferentes, mais núcleos físicos disponíveis, cada processo usando a memória do seu próprio nó. Mais \emph{threads} por processo significam menos mensagens, porém mais sincronização entre \emph{threads} e mais núcleos disputando a memória de um mesmo nó. Para $p = \Valor{hibridoP}$, a diferença medida entre a melhor e a pior divisão foi de $\Valor{hibridoDiferencaPct}\pct$, a favor de \Valor{hibridoMelhorNP} processo(s) com \Valor{hibridoMelhorT} \emph{thread}(s).

O padrão se repete nos três grupos da Tabela~\ref{tab:hibrido}: para o mesmo total de núcleos, a divisão com mais processos é sempre a mais rápida. Essa comparação precisa ser lida com cuidado, porque as divisões não usam a mesma máquina. Como cada processo fica em um nó, quatro processos de duas \emph{threads} usam quatro nós, e dois processos de quatro \emph{threads} usam dois. O que os números mostram é que usar um nó a mais compensa muito mais do que usar mais núcleos do mesmo nó, e que a comunicação entre nós custa pouco nesta aplicação.

Para este \emph{cluster}, a recomendação é usar um processo por nó. Quanto às \emph{threads}, o menor tempo foi obtido com quatro por processo, que é o número de núcleos físicos do nó e a distribuição prevista no mapeamento. O ganho da terceira e da quarta \emph{thread}, porém, é pequeno: na carga \Valor{cargaRef}, quatro processos levam $0{,}76$\,s com duas \emph{threads} e $0{,}71$\,s com quatro. Quando o objetivo é o menor tempo, vale usar as quatro. Quando o objetivo é aproveitar bem os núcleos ocupados, duas \emph{threads} por processo entregam quase o mesmo tempo com metade dos núcleos.

\subsection{Fração serial e limite do \emph{speedup}}

A métrica de Karp--Flatt da Tabela~\ref{tab:escalabilidade} separa duas situações. Se $e(p)$ permanece aproximadamente constante com o aumento de $p$, o limite do \emph{speedup} vem de uma fração serial fixa do algoritmo (Lei de Amdahl). Se $e(p)$ cresce com $p$, o limite vem do \emph{overhead} de paralelização, que aumenta com o número de núcleos, como a comunicação e a sincronização.

Dentro do trecho cronometrado, a parte que não é dividida entre os núcleos é pequena: a soma dos contadores, a troca dos ponteiros e a atualização do pico, feitas por uma \emph{thread} em tempo constante por passo.

Na Tabela~\ref{tab:escalabilidade}, $e(p)$ cresce com $p$: vale $-0{,}072$ com dois núcleos, $-0{,}018$ com quatro, $0{,}006$ com oito e $0{,}063$ com dezesseis. Os valores negativos aparecem porque a versão paralela com um núcleo é mais rápida que a sequencial usada no cálculo. Como $e(p)$ não é constante, o limite do \emph{speedup} não vem de uma fração serial fixa do algoritmo, e sim de um custo que aumenta com o número de núcleos. O maior salto ocorre de oito para dezesseis núcleos, que é justamente a passagem de duas para quatro \emph{threads} por processo discutida acima. Por esse motivo, este relatório não usa a Lei de Amdahl para projetar um \emph{speedup} máximo: a projeção suporia uma fração serial constante, o que as medidas contradizem.

Quando o problema aumenta, o custo de comunicação por passo, proporcional a $C$, cresce mais devagar que o trabalho útil, proporcional a $L \cdot C/p$. Pela Lei de Gustafson, isso indica que matrizes maiores tendem a ter eficiência maior para um mesmo $p$. Nas medidas, isso aconteceu da carga pequena para a média, mas não da média para a grande, e matrizes maiores que a grande não foram testadas.

\subsection{Limitações}

Quatro limitações da solução foram identificadas. A geração da matriz e o \emph{checksum} são sequenciais na raiz, por exigência do enunciado, e não aparecem no tempo cronometrado, mas dominam o tempo total de execução quando o núcleo da simulação é acelerado. A raiz mantém a matriz global, o que limita o tamanho do problema à memória de um nó. Por fim, a divisão da matriz apenas por linhas limita o número útil de processos a uma fração de $L$, pois blocos com poucas linhas gastam proporcionalmente mais tempo trocando fronteiras do que calculando. Por último, dentro de um nó o ganho com \emph{threads} praticamente se esgota em duas, e a causa disso não foi determinada.

\section{Conclusão}
\label{sec:conclusao}

A versão híbrida em C/MPI/OpenMP produziu resultados \Valor{resultadoValidacao} aos da versão sequencial do primeiro trabalho em todas as \Valor{numExecucoesPar} execuções cronometradas e em toda a bateria de casos de borda, para todas as combinações de processos e \emph{threads} testadas, o que confirma que a divisão da matriz em blocos de linhas, com as linhas fantasmas, não altera o resultado do modelo (Seção~\ref{sec:validacao}).

O projeto pela metodologia PCAM levou a uma solução em que a comunicação por passo se resume a duas linhas por processo e a uma redução de seis inteiros, toda ela feita com primitivas coletivas: envio dos parâmetros, distribuição da matriz, troca de linhas entre processos vizinhos, redução global e junção do resultado. Dentro de cada processo, a região paralela persistente com \texttt{omp for}, \texttt{simd}, reduções e tarefas dispensou \texttt{critical} e \texttt{atomic}, e a separação das linhas de fronteira em tarefas permitiu calcular as linhas internas do bloco enquanto as mensagens são trocadas (Seções~\ref{sec:pcam} e~\ref{sec:impl}).

Na carga \Valor{cargaRef}, o melhor tempo foi obtido com \Valor{melhorNP} processos de \Valor{melhorT} \emph{thread}(s), com \emph{speedup} absoluto de $\Valor{melhorSpeedupAbs}\times$ e eficiência de $\Valor{melhorEficienciaAbs}\pct$ sobre \Valor{melhorP} núcleos. O ganho com o número de processos foi próximo do ideal, com cada processo em um nó, o que mostra que a comunicação por passo custa pouco. O ganho com \emph{threads} dentro de um nó foi próximo do ideal até duas e pequeno de duas para quatro, e é isso que limita a eficiência da melhor configuração. Para um mesmo total de núcleos, a divisão com mais processos, e portanto mais nós, foi sempre a mais rápida.

\appendix

\section{Guia de compilação e execução}
\label{sec:apendice-execucao}

Esta seção resume o necessário para compilar e executar a versão paralela a partir dos arquivos entregues (\texttt{fire\_mpi\_omp.c} e \texttt{Makefile}), seguindo a forma de invocação exigida no enunciado.

\begin{lstlisting}[language=make,caption={Compilação.},basicstyle=\ttfamily\footnotesize]
make all      # compila a versao paralela: gera o binario fire_mpi_omp
make clean    # remove o binario gerado, se for preciso recompilar do zero
\end{lstlisting}

O alvo \texttt{all} chama o \texttt{mpicc} com as \emph{flags} \texttt{-Wall -Wextra -O2 -fopenmp -std=c99}, portanto não é necessário passar nenhuma \emph{flag} manualmente.

\begin{lstlisting}[language=bash,caption={Execução, no formato exigido pelo enunciado.},basicstyle=\ttfamily\footnotesize]
mpirun --hostfile hosts.txt -np 4 ./fire_mpi_omp entrada.txt
\end{lstlisting}

O programa recebe apenas o caminho do arquivo de entrada como argumento e o processo raiz escreve as doze linhas de saída em \texttt{stdout}. O número de processos é o informado ao \texttt{mpirun}, e o número de \emph{threads} $T$ de cada processo é lido da primeira linha do arquivo de entrada, não de uma variável de ambiente ou de um argumento adicional. O arquivo de entrada precisa estar acessível apenas ao processo raiz.

Uma observação sobre o Open MPI: quando NPrc é 1 ou 2, o \texttt{mpirun} prende por padrão cada processo a um único núcleo, e as $T$ \emph{threads} do processo passam a disputar esse núcleo. O resultado continua correto, mas o tempo deixa de refletir o paralelismo das \emph{threads}. Nos experimentos deste relatório isso foi evitado com a opção \texttt{--bind-to none}.

\section{Tabelas das demais cargas}
\label{sec:apendice-tabelas}

\begin{table}[H]
\centering
\small
\InputPlot{tab_tempos_pequena.tex}
\caption{Tempo do trecho cronometrado (mediana, em segundos) em função de NPrc e de $T$, carga pequena.}
\end{table}

\begin{table}[H]
\centering
\small
\InputPlot{tab_speedup_pequena.tex}
\caption{\emph{Speedup} absoluto em função de NPrc e de $T$, carga pequena.}
\end{table}

\begin{table}[H]
\centering
\small
\InputPlot{tab_eficiencia_pequena.tex}
\caption{Eficiência absoluta em função de NPrc e de $T$, carga pequena.}
\end{table}

\begin{table}[H]
\centering
\small
\InputPlot{tab_tempos_media.tex}
\caption{Tempo do trecho cronometrado (mediana, em segundos) em função de NPrc e de $T$, carga média.}
\end{table}

\begin{table}[H]
\centering
\small
\InputPlot{tab_speedup_media.tex}
\caption{\emph{Speedup} absoluto em função de NPrc e de $T$, carga média.}
\end{table}

\begin{table}[H]
\centering
\small
\InputPlot{tab_eficiencia_media.tex}
\caption{Eficiência absoluta em função de NPrc e de $T$, carga média.}
\end{table}

\end{document}